\begin{statement}{E.51}{Proved}{The prime-seven five-strand refinement preserves the full preceding family}
Let $\rho_7:\widehat B\twoheadrightarrow G_7$ be the actual source
obtained by inducing all eight evaluations of the four chain
transvections in $\operatorname{Sp}_4(\mathbb F_7)$ through all six
cosets of $B/F$, with the same words as E.50.  Its exact finite
images are
\[
\begin{gathered}
 |G_7|=135\,531\,648,\qquad
 |\rho_7(F)|=|\rho_7(H)|=22\,588\,608,\\
 C_7:=\rho_7(F')=\rho_7(H')=\rho_7(\gamma_rH)
       \quad(r\ge2),\qquad |C_7|=115\,248,\\
 L_7:=C_7'=\rho_7(N')\simeq\operatorname{SL}_2(\mathbb F_7),\\
 C_7=L_7\times Q_7,\quad Q_7\simeq(\mathbb Z/7)^3,\\
 \rho_7(N)=L_7\times(\mathbb Z/7)^5.
\end{gathered}
\tag{E.51.1}
\]
The isomorphisms in this display are the specified matrix and
generator maps in the proof.

Write
\[
\begin{gathered}
 A_{14}=B/(F'F^{14}),\qquad A_2=B/(F'F^2),\\
 Q_{28}=\operatorname{SL}_2(\mathbb Z/28)/\{\pm I\}.
\end{gathered}
\tag{E.51.2}
\]
The signed modular map and the six-coset abelian map give the
actual quotient
\[
\begin{gathered}
 G_7/Q_7\simeq A_{14}\times_{A_2}Q_{28},\\
 |A_{14}|=1176,\quad |Q_{28}|=8064,\quad
 |G_7/Q_7|=395\,136.
\end{gathered}
\tag{E.51.3}
\]
The maximum common quotient of $G_7$ and the E.38 source $G_L$
through their specified maps from $B$ is $A_2$.

Let $G_{\rm old}$ be any cofinal finite common source of the
entire E.50 family that retains
$P_{28}=B/(\gamma_3(H)H^{28})$.  It includes the fixed
$D_{23}$, $G_L$, and prime-five symplectic source maps, together
with any prescribed finite levels of the boundary factors.
Then the maximum common quotient of $G_{\rm old}$ and $G_7$
is the displayed $A_{14}$, and the actual joint source is
\[
                   G_{\rm old}\times_{A_{14}}G_7.
\tag{E.51.4}
\]
All $432$ retained multiplier-one source points, with every
prescribed full two-adic class-two fibre of E.41 and full
three-adic class-three correction-image fibre of E.47, have
one simultaneous profinite commutator word satisfying this
new test.  It retains all the full pro-$p$ $H$ factors from
E.42, the complete E.50 prime-five test, and the E.46 Hurwitz
test.  In both new symplectic sources and in $G_L$, this word
has exactly the values of the same actual Galois element
$\sigma$ with $\chi(\sigma)=1$.  Its changed source map is
onto \textup{(E.51.4)} at every finite level.

Put $T_7=\operatorname{Sp}_4(\mathbb F_7)$ and let $J_7$ be
its actual first four-strand image.  Then
\[
 |T_7|=276\,595\,200,\qquad |J_7|=115\,248,\qquad J_7'=J_7.
\tag{E.51.5}
\]
With $J_5$ the corresponding order-$15000$ group in E.50,
the complete joint four-strand target for this family is
\[
       (T_{23}\times_{S_L}T_L)\times J_5\times J_7\times A_{96}.
\tag{E.51.6}
\]
The same word generates it and satisfies all its braid and
full-twist identities.  The corresponding full five-strand
target is
\[
          \operatorname{Sp}_4(\mathbb F_5)
               \times\operatorname{Sp}_4(\mathbb F_7)\times A_{96},
\tag{E.51.7}
\]
again with the actual simultaneous generator maps.  These
assertions retain the finite sources and the indicated
boundary inverse limits; they impose no additional identity
between the full profinite GT group and its Galois subgroup.
\end{statement}

\begin{proof}
All comparisons below retain the displayed maps from the
fixed source $B$.

\paragraph{The literal representation and residual maps.}
Use the alternating matrix $\Omega$ whose only nonzero
upper-diagonal entries are $\Omega_{i,i+1}=1$, and put
$U_i=I+e_i(e_i^t\Omega)$ over $\mathbb F_7$.
The program with prefix
\texttt{bc\_symplectic\_five\_p7\_model}
checks the symplectic identities, all Artin relations and the
eight specified pure-braid evaluation pairs.  The action on
the $2400$ nonzero vectors is faithful and has order
$7^4(7^2-1)(7^4-1)=276595200$.  It is therefore the whole
symplectic group.  The first three generators have image
order $115248$.  The periodic words $\delta_5$ and
$\epsilon_5$ have orders ten and eight, respectively, and
$\delta_5^5=\epsilon_5^4=-I$.

The induced source starts on $6\cdot8\cdot2400=115200$
points.  The source certificate retains explicit
intertwiners from all $592$ original orbits to ten distinct
$B$ sets, of total degree $3570$.  Thus its compression is
an isomorphism of the actual permutation image.  The
subgroup calculation evaluates all $97$ literal $N$
generators and all $95$ primitive columns imposing zero
integral $F$ abelianization.  It gives
\[
\begin{array}{c|r}
\text{actual image}&\text{order}\\\hline
G_7&135531648\\
\rho_7(F)=\rho_7(H)&22588608\\
\rho_7(N)=\rho_7(K_6)&5647152\\
\rho_7(F')=\rho_7(K_6\cap F')&115248\\
\rho_7(N')=C_7'&336.
\end{array}
\tag{E.51.8}
\]
It also checks directly that
$[\rho_7(H),C_7]=C_7$.  Consequently every lower-central
term from the second onward has this same image.

Here is the exact matrix comparison underlying the direct
products.  Decode the identity-coset third pentagon
coordinate through the saved intertwiners.  It preserves
$\langle e_2,e_3\rangle$; transpose its block to account for
the right-permutation product.  This gives a homomorphism
$\theta:\rho_7(F)\to\operatorname{SL}_2(\mathbb F_7)$ with
\[
 \theta(x)=\begin{pmatrix}1&0\\2&1\end{pmatrix},
 \qquad\theta(y)=\begin{pmatrix}1&5\\0&1\end{pmatrix}.
\tag{E.51.9}
\]
Its restriction to $L_7$ takes all $336$ values exactly
once.  For every original generator $g$ of $\rho_7(N)$ or
$C_7$, let $\ell_g$ be the unique element of $L_7$ with
$\theta(\ell_g)=\theta(g)$, and put
\[
                         v_g=g\ell_g^{-1}.
\tag{E.51.10}
\]
The literal permutations verify that all these $v_g$
commute with one another and with $L_7$, and have seventh
power one.  They generate groups of orders $7^5$ and $7^3$,
respectively.  Since $\theta$ is injective on $L_7$ and
kills these groups, their intersections with $L_7$ are
trivial.  The displayed generator decompositions prove
the actual direct products in \textup{(E.51.1)}.

Define the ordinary modular map $q_7$ using
\[
 s\longmapsto\begin{bmatrix}0&1\\-1&0\end{bmatrix},\qquad
 r\longmapsto\begin{bmatrix}0&-1\\1&-1\end{bmatrix}
             \quad\text{in }\operatorname{PSL}_2(\mathbb F_7).
\tag{E.51.11}
\]
For $P=\left(\begin{smallmatrix}0&1\\1&0\end{smallmatrix}\right)$,
all $672$ conjugations of the $336$ elements of $L_7$ by
$s,r$ satisfy
\[
 \theta(a^{-1}\ell a)
       =(Pq_7(a)P)^{-1}\theta(\ell)(Pq_7(a)P).
\tag{E.51.12}
\]
The formula is independent of the sign of a representative
of $q_7(a)$.  In particular the induced action on the
inner automorphism group is exactly the specified
projective modular map, with this coordinate swap.

The certificate enumerates every normal closure in
$\operatorname{SL}_2(\mathbb F_7)$: the identity gives
order one, $-I$ gives order two, and each of the other
$334$ elements gives the whole order $336$ group.
The literal derived subgroup is the whole group.
Thus its nontrivial perfect quotients are precisely
itself and $\operatorname{PSL}_2(\mathbb F_7)$.
It follows as well that
\[
 \text{every nontrivial quotient of }C_7
                    \text{ has order divisible by }7.
\tag{E.51.13}
\]
Indeed an image of $Q_7$ has seventh-power order, and
every nontrivial image of $L_7$ has order divisible by
seven.  These images generate any quotient of their
direct product.

\paragraph{Both signed modular quotients and their actual old-source comparison.}
The six-coset Schreier exponent sums modulo fourteen
give $A_{14}$, with kernel $F'F^{14}$ and order $1176$.
The faithful joint image with $G_7$ still has order
$135531648$, proving the quotient map and its kernel
$C_7$.  Reducing the same coordinates modulo two gives
the exact $A_2$ map.

Enumerate $Q_{28}$ using the matrices in
\textup{(E.51.11)} modulo twenty-eight, retaining the
global sign.  It has $8064$ elements.  Its reduction
modulo four is identified with $A_2$ on all $24$ values,
using the same dictionary as E.38 and E.50.  An element
in this kernel has a unique sign representative equal
to $I$ modulo four; reduction of that representative
modulo seven identifies the kernel with all of
$\operatorname{SL}_2(\mathbb F_7)$.

The six-coset representation induced from the two
signed $\operatorname{SL}_2(\mathbb F_7)$ matrices
$\left(\begin{smallmatrix}1&2\\0&1\end{smallmatrix}\right)$,
$\left(\begin{smallmatrix}1&0\\-2&1\end{smallmatrix}\right)$
is compared with every element and all $16128$
generator edges of $Q_{28}$.  Thus it is precisely
this signed quotient.  All $395136$ matching pairs
in $A_{14}\times_{A_2}Q_{28}$ are generated by the
literal $s,r$.  Their joint map with $G_7$ retains
its original order; the remaining kernel has
order $343$ and is $Q_7$ from
\textup{(E.51.10)}.  This proves
\textup{(E.51.3)} with its specified maps.

The old-source modular comparison is different from
the prime-five computation and is checked directly.
Pair the faithful $G_A$ source with the action of
\textup{(E.51.11)} on $\mathbb P^1(\mathbb F_7)$.
The exact joint order is $168|G_A|$, so
\[
       \operatorname{im}(\rho_A,q_7)
           =G_A\times\operatorname{PSL}_2(\mathbb F_7).
\tag{E.51.14}
\]
The E.38 map $G_L\to G_A$ has elementary two-group
kernel.  Therefore
$q_7(\ker\rho_A)/q_7(\ker\rho_L)$ is a two-group.
Its numerator is the perfect
$\operatorname{PSL}_2(\mathbb F_7)$ by
\textup{(E.51.14)}, so the quotient is trivial.
Thus $q_7(\ker\rho_L)$ is the whole projective group.

Because the actual $Q_4=A_2$ map factors through $G_L$,
the image $J$ of $\ker\rho_L$ in $Q_{28}$ is contained
in its signed kernel $L=\operatorname{SL}_2(\mathbb F_7)$
and maps onto $L/Z(L)$.  Hence $JZ(L)=L$, and
\[
             L=[L,L]=[JZ(L),JZ(L)]=[J,J]\subset J.
\tag{E.51.15}
\]
This proves the actual joint image
$G_L\times_{A_2}Q_{28}$, including its central sign.

\paragraph{The complete correction preserving E.50 and every boundary factor.}
Put
$R_{\rm nil}=\bigcap_{r\ge1}\gamma_r(\widehat H)$,
and use $C_L$ for the perfect E.43 residual in $G_L$.
For any finite image of $\widehat H$, its lower-central
residual is the image of $R_{\rm nil}$: surjectivity
at each lower-central term and nested compact fibres
prove the reverse inclusion as well as the immediate
forward one.  E.43 and \textup{(E.51.1)} consequently give
\[
                   \rho_L(R_{\rm nil})=C_L,
          \qquad \rho_7(R_{\rm nil})=C_7.
\tag{E.51.16}
\]
For the first equality, the quotient of the $H$ image
in $G_L$ by $C_L$ is the nilpotent group of E.43;
the closed derived residual already maps onto $C_L$.

The joint image in \textup{(E.51.16)} is the full product.
To prove this, any nontrivial Goursat quotient is
perfect, since $C_L$ is perfect.  A perfect quotient
of $C_7=L_7\times Q_7$ is $L_7$ or its projective
quotient.  Indeed the image of $L_7$ must be the whole
perfect quotient; its centre has order at most two,
so the central seventh-power group $Q_7$ maps to one.
Both Goursat kernels are stable under the full $B$
action, because $R_{\rm nil}$ is characteristic in
the normal subgroup $\widehat H$.  Formula
\textup{(E.51.12)} would therefore make its projective
inner action factor through $G_L$.  This would factor
the specified $q_7$ through $G_L$, contrary to
\textup{(E.51.14)--(E.51.15)}.  Hence
\[
             \operatorname{im}_{G_L\times G_7}(R_{\rm nil})
                              =C_L\times C_7.
\tag{E.51.17}
\]

In particular $K=R_{\rm nil}\cap\ker\rho_L$ maps onto
$C_7$.  The image of $K$ in the actual
$D_{23}\times G_{5,\mathrm{sp}}$ has order prime to
seven.  By \textup{(E.51.13)}, its joint image with
$C_7$ is the full product.  Intersecting with this
finite kernel therefore retains the surjection
\[
 R_{\rm nil}\cap\ker\rho_L
      \cap\ker(\rho_{D_{23}},\rho_{5,\mathrm{sp}})
                      \twoheadrightarrow C_7.
\tag{E.51.18}
\]
This is the required explicit comparison with the
E.50 kernel.  Every finite nilpotent $H$ image kills
$R_{\rm nil}$, so the same closed subgroup kills
all finite boundary factors and their full indicated
inverse limits.  No separate correction for a
boundary level is needed.

The pure $C_7$ in \textup{(E.51.17)} forces every
common quotient of $G_L,G_7$ through $A_{14}$.
This is solvable, whereas the maximal solvable
quotient of $G_L$ has order $43200$ and no factor
seven.  Its common quotient kills the $(\mathbb Z/7)^2$
kernel of $A_{14}\to A_2$.  The specified $A_2$
maps attain the resulting quotient, proving its
maximality.

For $G_{\rm old}$, the actual map $P_{28}\to A_{14}$
follows from $\gamma_3H\subset F'$ and
$H^{28}\subset F^{14}$.  The $P_4$ and full pro-seven
$H$ factors already supply $P_{28}$ by their actual
coprime geometric-image product.  More explicitly,
the class-two exponent-$28$ quotient of $H$ is the
direct product of its exponent-four and exponent-seven
Sylow quotients; their maps are the original source
reductions, and the inverse uses the two coprime
power idempotents.  Thus requiring $P_{28}$
only passes to a cofinal subsystem.  The pure $C_7$
in \textup{(E.51.18)} forces every common quotient
through $A_{14}$, and $A_{14}$ is a specified common
quotient.  Goursat gives exactly
\textup{(E.51.4)}.

\paragraph{One Galois comparison and full joint generation.}
Choose the E.50 word $f_0$ for the prescribed old
point and boundary parameters.  Retain its actual
Galois element $\sigma$ with
$\chi(\sigma)=1\in\widehat{\mathbb Z}^{\times}$.
In particular $\chi(\sigma)\equiv1\pmod{280}$,
which retains the orders fourteen of the new
source generators, seven of the transvections,
and ten and eight of the periodic words.
Since both words belong to $\widehat F'$, their
new images differ by
\[
          c=\rho_7(f_0)^{-1}\rho_7(f_\sigma)\in C_7.
\tag{E.51.19}
\]
Lift $c$ through \textup{(E.51.18)} and put $f=f_0k$.
This is a single word in $\widehat F'$.  It retains
every old value and satisfies
$\rho_7(f)=\rho_7(f_\sigma)$.  All $48$ induced
evaluation values agree with this same Galois word.
The new relations, periodic conjugacies and
individual generation are consequently the actual
Galois braid maps followed by the specified finite
quotient maps.  This argument does not require
an arbitrary finite shadow to preserve its source
kernel.

For joint generation, the corrected source maps
to $G_L$ and $G_7$ are both precomposed with the
same actual Galois automorphism of $\widehat B$.
It preserves $\widehat F$, its characteristic
subgroup $\widehat H$, and $R_{\rm nil}$.
Thus their joint residual image remains exactly
\textup{(E.51.17)}.  Intersecting with the corrected
$D_{23}$ and prime-five symplectic kernels uses
the same prime-to-seven argument in
\textup{(E.51.18)}.  The corrected boundary maps
still kill $R_{\rm nil}$, since their $H$ images
are nilpotent at each finite level.  The corrected
joint image therefore contains the full pure
$C_7$.  Its projection onto $G_{\rm old}$ is onto
by E.50, and its maps to $A_{14}$ agree because
they are the same literal source formulas.
The quotient of \textup{(E.51.4)} by pure $C_7$
is $G_{\rm old}$, proving surjectivity onto the
full fibre product.  The actual $H$ restrictions
are onto as well: the whole maps are onto and
their retained $B/H$ map is the same automorphism.
Nested compact fibres give the asserted one-word
comparison through every boundary level.

The image parameters do not specify a unique
profinite word.  If $\mathcal K$ denotes the closed
correcting subgroup on the left of
\textup{(E.51.18)}, the corrections within this subgroup
for the fixed value $c$ are exactly
$k(\mathcal K\cap\ker\rho_7)$.  If $\rho_{\rm old}$
denotes the full joint old map on $\widehat F'$, the
set of all words with the prescribed old value and
the prescribed new Galois value is exactly
\[
       f\,\ker\bigl((\rho_{\rm old},\rho_7)
                                      |_{\widehat F'}\bigr).
\tag{E.51.20}
\]
Both assertions follow by multiplying two preimages
under the displayed continuous homomorphisms.

\paragraph{The joint braid targets.}
The literal commutators $[U_1,U_2]$,
$[U_2,U_3]$, and $U_1[U_1,U_2]U_1^{-1}$ generate
all $115248$ matrices of $J_7$; the intermediate
orders are $2744$ and $115248$.  This proves
$J_7'=J_7$ without a structural identification
of its unipotent radical.

The complete lifted Schreier calculation pairs
the three actual Ferrari generators in $T_L$
with $U_1,U_2,U_3$.  It retains all $57$ ambient
binary coordinates and recovers the full
$55$-dimensional central kernel above the
actual permutation product.  Its exact order
is $|T_L|\,115248$, proving the actual target
$T_L\times J_7$.  A separate literal paired
permutation calculation gives $J_5\times J_7$
and, on all four generators,
$\operatorname{Sp}_4(\mathbb F_5)\times T_7$.

These comparisons also give the joint
$T_L\times J_5\times J_7$.  In a nontrivial
Goursat quotient of $T_L\times J_5$ and $J_7$,
quotient by either of the two commuting factor
images.  A nontrivial result contradicts one
of the two checked pair products.  If both
results are trivial, those commuting images
both equal the quotient, so it is abelian.
It is also perfect, as a quotient of $J_7$,
and hence is trivial after all.

The kernel of
$(T_{23}\times_{S_L}T_L)\times J_5
 \to T_L\times J_5$ is solvable.  A nontrivial
perfect Goursat quotient with $J_7$ cannot be
the image of this solvable kernel; quotienting
by that image would give a nontrivial common
quotient with $T_L\times J_5$, already excluded.
This proves the complete product before the
$A_{96}$ factor in \textup{(E.51.6)}.
The simple group $A_{96}$ has order divisible
by eleven, while the other factors do not.
The actual Hurwitz projection and Goursat
therefore give its full product as well.

All the $T_L,J_5,J_7,A_{96}$ components of
the corrected word are the same Galois
precomposition, so their simultaneous images
are the actual products just proved.  The
solvable-kernel argument retains the old
$T_{23}$ component and proves corrected
generation in \textup{(E.51.6)}.  The paired
five-strand calculation and the same
$A_{96}$ argument prove
\textup{(E.51.7)} with its corrected maps.

The new programs have the shared prefix
\texttt{bc\_symplectic\_five\_p7\_}.
Their suffixes \texttt{model}, \texttt{source},
\texttt{subgroups}, \texttt{residual},
\texttt{modular\_probe}, and \texttt{common}
retain all source maps and signed comparisons;
\texttt{target}, \texttt{target\_perfect}, and
\texttt{joint\_symplectic} retain the full
target comparisons.  Each proof assertion
above uses these actual finite maps and
their stated compact inverse-limit maps.
\end{proof}
